% ============================================
% MEM-07
% Gerado de molambudos.md — build_miolo.py
% ============================================
\section*{\textbf{MEM-07} --- A Mãe Levada}
\markboth{MEM-07 --- A Mãe Levada}{}

\noindent \textit{Soldados, caminhão, silêncio}

\noindent \textbf{1915 --- O Curral do Governo}

\vspace{0.5em}\noindent\rule{\textwidth}{0.5pt}\vspace{0.5em}

\noindent Depois da noite na vala, minha mãe ficou diferente.

\noindent Ela me olhava de um jeito que eu não conhecia. Não era o olhar de mãe --- era o olhar de alguém que viu um bicho e não sabe se o bicho é perigoso.

\noindent \textit{\textquotedblleft{}Teu olho\textquotedblright{}\textit{, ela repetia. }\textquotedblleft{}Teu olho, Joaquim.\textquotedblright{}}

\noindent O olho direito estava amarelo. Não amarelo claro. Amarelo de ovo. Amarelo de gema podre. Eu via o reflexo nos cacos de louça que a gente usava para beber água.

\noindent Eu dizia que não era nada.

\noindent Ela sabia que era mentira.

\vspace{0.5em}\noindent\rule{\textwidth}{0.5pt}\vspace{0.5em}

\noindent Os soldados chegaram numa tarde.

\noindent Eram quatro. Dois com espingardas. Um com chicote. Um com uma lista de nomes.

\noindent \textit{\textquotedblleft{}Mulheres\textquotedblright{}\textit{, gritou o do chicote. }\textquotedblleft{}Todas as mulheres entre quinze e quarenta anos. Apresentem-se na entrada.\textquotedblright{}}

\noindent Minha mãe me apertou contra o peito.

\noindent \textit{\textquotedblleft{}Fica quieto\textquotedblright{}\textit{, ela sussurrou. }\textquotedblleft{}Fica quieto que eles não me levam.\textquotedblright{}}

\noindent Mas levaram.

\noindent O soldado com a lista leu o nome dela. Ela não respondeu. Ele leu de novo. Ela não respondeu. Ele veio até nós.

\noindent \textit{\textquotedblleft{}Você é Maria do Rosário Correia?\textquotedblright{}}

\noindent Ela não respondeu.

\noindent \textit{\textquotedblleft{}Você tá na lista. Vem.\textquotedblright{}}

\noindent \textit{\textquotedblleft{}Meu filho\textquotedblright{}\textit{, ela disse. }\textquotedblleft{}Meu filho precisa de mim.\textquotedblright{}}

\noindent \textit{\textquotedblleft{}Teu filho vai ficar bem. Os meninos ficam. As mulheres vão pra outra cidade. Lá tem comida.\textquotedblright{}}

\noindent \textit{\textquotedblleft{}Mentira\textquotedblright{}\textit{, ela disse. }\textquotedblleft{}Eu sei o que vocês fazem com as mulheres.\textquotedblright{}}

\noindent O soldado não respondeu. Ele fez um sinal com a cabeça. Os outros dois vieram. Um segurou o braço esquerdo dela. O outro segurou o direito.

\noindent \textit{\textquotedblleft{}Mãe!\textquotedblright{}}, eu gritei.

\noindent \textit{\textquotedblleft{}Fica, Joaquim! Fica!\textquotedblright{}}, ela gritou de volta.

\noindent Ela me olhou. Os olhos dela estavam cheios d'água. Mas não era choro. Era raiva. Uma raiva que eu nunca tinha visto.

\noindent \textit{\textquotedblleft{}Não morre\textquotedblright{}\textit{, ela disse. }\textquotedblleft{}Não morre, Joaquim. Escuta a tua mãe: não morre.\textquotedblright{}}

\noindent Os soldados arrastaram ela.

\noindent Eu corri atrás.

\noindent O do chicote me empurrou. Eu caí no chão. Levantei e corri de novo. Ele me empurrou de novo. Dessa vez eu não levantei.

\noindent Fiquei no chão, olhando.

\noindent O caminhão partiu.

\noindent Minha mãe estava na carroceria, entre outras mulheres. Ela me olhava. Eu olhava ela.

\noindent O caminhão ficou pequeno.

\noindent Depois sumiu.

\noindent Eu nunca mais vi minha mãe.\NE{A \textquotedblleft{}realocação\textquotedblright{} e o destino de Maria são elementos ficcionais. A violência sofrida por mulheres retirantes exige documentação histórica específica; como Joaquim não existiu, sua narrativa não constitui testemunho nem confirma um arquivo oficial. O conflito entre a promessa de comida e o medo da separação move este fragmento apenas dentro da ficção.}

\vspace{0.5em}\noindent\rule{\textwidth}{0.5pt}\vspace{0.5em}

\noindent O velho Anselmo, que não tinha morrido ainda naquela época, me disse depois:

\noindent \textit{\textquotedblleft{}Elas vão pros bordéis do litoral. O governo chama de 'realocação'. As que não querem ir, vão pro mato com um tiro na nuca.\textquotedblright{}}

\noindent \textit{\textquotedblleft{}E minha mãe?\textquotedblright{}}

\noindent \textit{\textquotedblleft{}Tua mãe foi.\textquotedblright{}}

\noindent \textit{\textquotedblleft{}Foi pra onde?\textquotedblright{}}

\noindent Anselmo cuspiu no chão.

\noindent \textit{\textquotedblleft{}Foi viver. Ou morrer. No fundo é a mesma coisa.\textquotedblright{}}

\noindent Eu deitei no chão do curral. Olhei pro céu. O céu estava limpo. Sem nuvem. Sem chuva. Sem nada.

\noindent Só azul.

\noindent Um azul que doía.

\noindent Eu senti a criatura se mexer dentro de mim.

\noindent \textit{\textquotedblleft{}Calma\textquotedblright{}\textit{, ela disse. }\textquotedblleft{}Um dia tu vai entender.\textquotedblright{}}

\noindent \textit{\textquotedblleft{}Entender o quê?\textquotedblright{}}

\noindent \textit{\textquotedblleft{}Que o mundo come os filhos. E os filhos aprendem a comer de volta.\textquotedblright{}}

\noindent Eu fechei os olhos.

\noindent O olho amarelo não fechou.

\vspace{0.5em}
\noindent\textit{\textbf{↪ Rotas:} \rota{MEM-06} \rotasep \rota{MEM-08} \rotasep \rota{DOC-05} \rotasep \rota{DOC-17}}